% 使用 XeLaTeX 编译；后续按板书照片顺序追加内容。
\documentclass[UTF8,a4paper,12pt]{ctexart}
\usepackage[margin=2.5cm]{geometry}
\usepackage{amsmath,amssymb,amsthm}
\usepackage{graphicx}
\usepackage[colorlinks=true,linkcolor=blue]{hyperref}
\newtheorem{proposition}{命题}

\title{复变函数课堂笔记}
\author{}
\date{}

\begin{document}
\maketitle

% 2026-10-19，第 1 张板书。
\setcounter{section}{4}
\section{留数（Residues）}

\begin{proposition}[解析函数零点的孤立性]
若 $f(z)$ 在 $z_0$ 点解析，$f(z_0)=0$，且 $f(z)\not\equiv 0$，
则存在 $\delta>0$，使得在 $|z-z_0|<\delta$ 内有
\[
  f(z)=(z-z_0)^n\varphi(z),
\]
这里 $n\in\mathbb{N}^{+}$，$\varphi(z_0)\ne 0$，且
$\varphi(z)$ 在 $|z-z_0|<\delta$ 内处处可导（解析）。
\end{proposition}

\paragraph{证明}
因 $f(z)$ 在 $z_0$ 解析，故存在 $\delta>0$，当 $|z-z_0|<\delta$ 时，有
\begin{align*}
  f(z)
  &=\sum_{k=0}^{+\infty}\frac{f^{(k)}(z_0)}{k!}(z-z_0)^k \\
  &=f(z_0)+\frac{f'(z_0)}{1!}(z-z_0)+\cdots
    +\frac{f^{(n)}(z_0)}{n!}(z-z_0)^n+\cdots.
\end{align*}
令
\[
  n=\min\bigl\{k: f^{(k)}(z_0)\ne 0\bigr\},
\]
则
\[
  f(z_0)=f'(z_0)=\cdots=f^{(n-1)}(z_0)=0.
\]
于是
\begin{align*}
  f(z)
  &=\frac{f^{(n)}(z_0)}{n!}(z-z_0)^n
    +\frac{f^{(n+1)}(z_0)}{(n+1)!}(z-z_0)^{n+1}+\cdots \\
  &=(z-z_0)^n\underbrace{\left[
      \frac{f^{(n)}(z_0)}{n!}
      +\frac{f^{(n+1)}(z_0)}{(n+1)!}(z-z_0)+\cdots
    \right]}_{\displaystyle\varphi(z)} \\
  &=(z-z_0)^n\varphi(z).
\end{align*}
% 第 1 张板书到此；后续证明按新照片继续追加。

\end{document}
